% Shared typography and code-listing setup for both book editions.
% Language selection and PDF metadata stay in each edition's main.tex.
\usepackage{lmodern}
\usepackage{microtype}
\usepackage[a4paper,margin=27mm]{geometry}
\usepackage{parskip}
\usepackage{needspace}
\usepackage{enumitem}
\usepackage{booktabs}
\usepackage{longtable}
\usepackage{tabularx}
\usepackage{xcolor}
\usepackage{listings}
\usepackage{hyperref}
\usepackage{bookmark}
\usepackage{fancyhdr}
\usepackage{upquote}
\usepackage{textcomp}

\setlist{nosep}
\pagestyle{fancy}
\fancyhf{}
\fancyhead[LE,RO]{\thepage}
\fancyhead[LO]{\nouppercase{\rightmark}}
\fancyhead[RE]{\nouppercase{\leftmark}}
\setlength{\headheight}{14pt}

% Ragged-right paragraph columns avoid stretched prose and ugly hyphenation in
% narrow reference tables while retaining explicit width control.
\newcolumntype{L}[1]{>{\raggedright\arraybackslash}p{#1}}
\newcolumntype{Y}{>{\raggedright\arraybackslash}X}
\setlength{\LTpre}{0.6em}
\setlength{\LTpost}{0.8em}

\makeatletter
\renewcommand*\l@chapter[2]{%
  \ifnum \c@tocdepth >\m@ne
    \addpenalty{-\@highpenalty}%
    \vskip 1.0em \@plus\p@
    \setlength\@tempdima{2.2em}%
    \begingroup
      \parindent \z@ \rightskip \@pnumwidth
      \parfillskip -\@pnumwidth
      \leavevmode \bfseries
      \advance\leftskip\@tempdima
      \hskip -\leftskip
      #1\nobreak\hfil \nobreak\hb@xt@\@pnumwidth{\hss #2}\par
      \penalty\@highpenalty
    \endgroup
  \fi}
\renewcommand*\l@section{\@dottedtocline{1}{1.5em}{4.0em}}
\renewcommand*\l@subsection{\@dottedtocline{2}{3.8em}{5.2em}}
\makeatother

\lstdefinelanguage{Velran}{
  morekeywords={model,object,query,fn,page,action,component,layout,form,route,GET,POST,auth,role,user,mfa,cache,public,ttl,transaction,sql,return,let,mut,mod,validate,range,length,same,pattern,enum,canonical,slug,from,changed,flash,invalidate,if,else,while,match,audit},
  sensitive=true,
  morecomment=[l]{//},
  morestring=[b]"
}
\lstset{
  basicstyle=\ttfamily\small,
  columns=fullflexible,
  keepspaces=true,
  breaklines=true,
  frame=single,
  showstringspaces=false,
  upquote=true,
  tabsize=4,
  aboveskip=0.8em,
  belowskip=0.8em
}
\newcommand{\code}[1]{\texttt{\detokenize{#1}}}
\newcommand{\velran}{Velran}
\newcommand{\callable}[1]{\texttt{\#[#1] fn}}
\newcommand{\invariantblock}[2]{%
  \Needspace{5\baselineskip}%
  \begin{quote}
  \noindent\textbf{#1}\par
  #2
  \end{quote}%
}
\newcommand{\diagnosticblock}[4]{%
  \Needspace{8\baselineskip}%
  \begin{quote}
  \noindent\textbf{#1: \code{#2}}\par
  \textbf{Meaning:} #3\par
  \textbf{Fix:} #4
  \end{quote}%
}
\newcommand{\diagnosticblockhu}[4]{%
  \Needspace{8\baselineskip}%
  \begin{quote}
  \noindent\textbf{#1: \code{#2}}\par
  \textbf{Jelentés:} #3\par
  \textbf{Javítás:} #4
  \end{quote}%
}
\newcommand{\rustbridgeblock}[3]{%
  \Needspace{8\baselineskip}%
  \begin{quote}
  \noindent\textbf{Rust $\rightarrow$ Velran}\par
  \textbf{Reuse from Rust:} #1\par
  \textbf{Velran adds:} #2\par
  \textbf{Do not relearn:} #3
  \end{quote}%
}
\newcommand{\rustbridgeblockhu}[3]{%
  \Needspace{8\baselineskip}%
  \begin{quote}
  \noindent\textbf{Rust $\rightarrow$ Velran}\par
  \textbf{Hozd Rustból:} #1\par
  \textbf{A Velran hozzáadja:} #2\par
  \textbf{Nem kell újratanulnod:} #3
  \end{quote}%
}
